\section{Problem Statement}\label{sec:problem}
Given a reference $q_d\in C^2([0,\infty);\mathbb R^n)$ with bounded
$q_d$, $\dot q_d$, and $\ddot q_d$, define
\begin{equation}
 e=q_m-q_d,\qquad \dot e=\dot q_m-\dot q_d,\qquad
 z=[e^\top,\dot e^\top]^\top.
 \label{eq:tracking_errors}
\end{equation}
The objective is to prescribe the total settling time through two design
parameters, $T_o$ and $h$:
\begin{equation}
 T=T_{\rm pre}=T_o+2h.
 \label{eq:prescribed_time_definition}
\end{equation}
\begin{problem}[Prescribed-time joint trajectory tracking]\label{prob:tracking}
Design $\tau(t)$ such that $e(t)=\dot e(t)=0$ for all $t\geq T$ under
ideal continuous-time exact compensation. Under nonidealities, establish an
explicit ultimate bound in terms of the residual acceleration disturbance.
\end{problem}

\subsection{Basic Assumptions}
\begin{assumption}[Equivalent inertia]\label{as:inertia}
There is $\underline m>0$ such that
$M_e(q_m)\succeq\underline m I_n$ on the considered workspace. The nominal
$M_e$ and $C_e$ used by the controller are available.
\end{assumption}
\begin{assumption}[Reference and measurements]\label{as:reference}
$q_d\in C^2$, and $q_d$, $\dot q_d$, and $\ddot q_d$ are bounded. The states
$q_m$ and $\dot q_m$ are measured and all signals used by the controller are
well defined on the considered interval.
\end{assumption}
\begin{assumption}[Delay history]\label{as:delay_history}
The function $z(t)$ is available on $[T_o-h,T_o]$ and the delayed signal is
implemented from a bounded, piecewise-continuous buffer. This history need not
satisfy the nominal model before $T_o$.
\end{assumption}
\begin{assumption}[Disturbance regularity]\label{as:disturbance}
The lumped disturbance is such that $\Delta_a=M_e^{-1}d$ is absolutely
continuous and there exist finite constants $\bar\Delta_a,L_d$ satisfying
\begin{equation}
 \|\Delta_a(t)\|\leq\bar\Delta_a,\qquad
 \|\dot\Delta_a(t)\|\leq L_d
\end{equation}
for almost every $t$. The observer gains are selected using the hypotheses of
the observer theorem stated in Section~\ref{sec:controller}.
\end{assumption}
\begin{assumption}[Ideal implementation]\label{as:ideal}
The exact prescribed-time result refers to continuous-time implementation,
absence of actuator saturation and measurement noise, exact model matching,
and an observer satisfying Assumption~\ref{as:observer_interface}. With
sampling, saturation, or model error, only the practical bound in
Section~\ref{sec:practical} is claimed.
\end{assumption}

\subsection{Delay-system preliminaries}
For a retarded linear delay system
\begin{equation}
 \dot x(t)=F(t)x(t)+G(t)x(t-h),
 \label{eq:abstract_delay_system}
\end{equation}
with a bounded initial history on an interval of length $h$, let $\Psi(t,s)$
denote its causal state-transition operator. If $F$ and $G$ are $2h$-periodic,
then the sampled state at period boundaries has a linear map
\begin{equation}
 x(2(k+1)h)=\Delta(h)x(2kh),\qquad k=0,1,\ldots,
 \label{eq:period_map_abstract}
\end{equation}
whenever the history has been included in the state representation. For the
PDF design below, $G(t)=0$ on the first half-period, so the delayed signal on
the second half-period is explicitly determined by the first-half trajectory
and the map can be calculated by variation of constants.

\begin{definition}[Single-period nullification]\label{def:single_period_nullification}
A PDF system has single-period nullification if its one-period map satisfies
$\Delta(h)=0$. In that case the state at the end of every complete period is
zero under ideal homogeneous dynamics, independently of the state entering the
period.
\end{definition}
\begin{definition}[Nilpotent nullification]\label{def:nilpotent_nullification}
If $\Delta(h)$ is nilpotent of index $\nu_0$, then the sampled state is zero
at or before $2\nu_0h$ under ideal homogeneous dynamics.
\end{definition}
